%! Author = Kevin Lin
%! Date = 10/8/2026

% Preamble
\documentclass[11pt,a4paper,margin=1in]{article}

% Packages
\usepackage{amsmath}
\usepackage{amssymb}
\usepackage{enumerate}
\usepackage{pdfpages}

\title{HW 2}
\author{Kevin Lin}
\date{10/8/2026}

% Document
\begin{document}
\maketitle

\section{}
\begin{enumerate}[(a)]
    \item Given $Z_1, \dots, Z_n$ are i.i.d. Bernoulli$(p)$ for $p \leq 1/2$
    and $t \in (0, 1 - p)$, we can compare the upper bound
    \[
        P\left[\frac{1}{n} \sum_{i=1}^{n} Z_i \geq t + p\right]
    \]
    between Hoeffding's inequality, Bernstein's inequality, and Bennett's inequality.
    First, define $X_i = Z_i - p$, consequently, $E[Z_i] = p, Var(Z_i) = p(1 - p)$,
    and thus $E[X_i] = 0, Var(X_i) = p(1 - p)$. By this same definition, we can derive
    that 
    \[
        P\left(\frac{1}{n} \sum_{i=1}^n Z_i \geq t + p \right) = P\left(\sum_{i=1}^n X_i \geq nt\right)
    \]
    Now for each inequality:
    \begin{itemize}
        \item \textbf{Hoeffding's inequality:} Hoeffding's
        inequality gives:
        \begin{gather*}
            P\left(\sum_{i=1}^{n} X_i \geq s) \leq \exp(-\frac{2s^2}{n})\right), \quad s = nt \therefore \\
            P\left(\frac{1}{n} \sum_{i=1}^{n} Z_i \geq t + p\right) \leq \exp(-2nt^2)
        \end{gather*}
        thus, $R_H = 2nt^2$.
        \item \textbf{Bernstein's inequality:} Using $X_i$, since $Z_i \in \{0, 1\}$,
        $X_i \in \{-p, 1 - p\}$. Because $p \leq 1/2$, $|X_i| \leq 1 - p$, and 
        that $\sum_{i=1}^n Var(X_i) = np(1 - p)$, Bernstein's inequality gives:
        \[
            P\left(\sum_{i=1}^{n} X_i \geq s\right) \leq \exp\left(-\frac{s^2}{2(v + bs/3)}\right)
        \]
        Again, let $s = nt$, with $v = np(1 - p)$ and $b = 1 - p$, we have:
        \begin{align*}
            P\left(\sum_{i=1}^{n} X_i \geq nt\right) &\leq \exp\left(-\frac{n^2 t^2}{2np(1 - p) + \frac{2}{3}(1 - p)nt}\right) \\
            &= \exp\left(-\frac{n^2t^2}{2n(1 - p)(p + t/3)}\right) \\
            &= \exp\left(-\frac{nt^2}{2(1 - p)(p + t/3)}\right)
        \end{align*}
        thus, $R_{\text{Bern}} = \frac{nt^2}{2(1 - p)(p + t/3)}$.
        \item \textbf{Bennett's inequality:} Using the same $X_i$ as above, Bennett's 
        inequality gives:
        \begin{align*}
            P\left(\sum_{i=1}^{n} X_i \geq s\right) &\leq \exp\left(-\frac{v}{b^2} h\left(\frac{bs}{v}\right)\right), \text{ using } s,v,b \therefore \\
            P\left(\sum_{i=1}^{n} X_i \geq nt\right) &\leq \exp\left(-\frac{np(1 - p)}{(1 - p)^2} h\left(\frac{(1 - p)nt}{np(1 - p)}\right)\right) \\
            &= \exp\left(-\frac{np}{1 - p} h\left(\frac{t}{p}\right)\right)
        \end{align*}
        thus, $R_{\text{Benn}} = \frac{np}{1 - p} h\left(\frac{t}{p}\right)$.
    \end{itemize}
    \item We can compute $R_{\text{Bern}}/R_H$ using $(a)$:
    \[
        \frac{R_{\text{Bern}}}{R_H} = \frac{\frac{nt^2}{2(1 - p)(p + t/3)}}{2nt^2} = \frac{1}{4(1 - p)(p + t/3)}
    \]
    We can see if Bernstein's bound improves over Hoeffding's more than a constant
    factor if $R_{\text{Bern}}/R_H \to \infty$ as $n \to \infty$. Since $p \leq \frac{1}{2}$,
    $\frac{1}{2} \leq 1 - p \leq 1$, so $1 - p$ is bounded away from 0. For the
    ratio to approach infinity, we need $p + t/3 \to 0$ as $n \to \infty$. This
    only occurs when $p \to 0$ and $t \to 0$ as $n \to \infty$. Thus, Bernstein's 
    bound improves over Hoeffding's more than a constant factor when $p \to 0$ and 
    $t \to 0$ as $n \to \infty$.
    \item $R_{\text{Benn}}/R_{\text{Bern}}$ depends on $(p, t)$ only through $x = t/p$:
    \begin{align*}
        \frac{R_{\text{Benn}}}{R_{\text{Bern}}} &= \frac{\frac{np}{1 - p} h\left(\frac{t}{p}\right)}{\frac{nt^2}{2(1 - p)(p + t/3)}} = \frac{2p(p + t/3)}{t^2} h\left(\frac{t}{p}\right) \\
        &= \frac{2p(p + px/3)}{p^2 x^2} h(x) = \frac{2(1 + x/3)}{x^2} h(x)
    \end{align*}
    Bennett's bound is never than Bernstein's since this ratio is always greater than 1:
    \begin{align*}
        \frac{2(1 + x/3)}{x^2} h(x) &\geq 1 \\
        h(x) &\geq \frac{x^2}{2(1 + x/3)} \\
        \text{Let } f(x) &= h(x) - \frac{x^2}{2(1 + x/3)} \\
        h(x) &= (1 + x)\log(1 + x) - x \therefore \\
        h'(x) &= \log(1 + x), h''(x) = \frac{1}{1 + x} \therefore \\
        f''(x) &= \frac{1}{1 + x} - \frac{1}{(1 + x/3)^3} \geq 0 
    \end{align*}
    For all $x \geq 0$, $(1 + x/3)^3 \geq 1 + x$, so $f''(x) \geq 0$, then $f'(x) \geq 0$ 
    also. With $f(0) = 0, f'(0) = 0$, we have $f(x) \geq 0$ for all $x \geq 0$. 
    Therefore, $h(x) \geq \frac{x^2}{2(1 + x/3)}$, and substituting this into the ratio gives:
    \[
        \frac{2(1 + x/3)}{x^2} h(x) \geq \frac{2(1 + x/3)}{x^2} \cdot \frac{x^2}{2(1 + x/3)} = 1
    \]
    thus, Bennett's bound is never worse than Bernstein's. \\
    As $x \to 0$, the taylor expansion of $h(x)$ gives $h(x) = \frac{x^2}{2} - \frac{x^3}{6} + O(x^4)$, so
    \[
        \frac{2(1 + x/3)}{x^2} h(x) = \frac{2(1 + x/3)}{x^2} \left(\frac{x^2}{2} - \frac{x^3}{6} + O(x^4)\right) = 1 + \frac{x^2}{18} + O(x)
    \]
    Thus, as $x \to 0$, the ratio approaches 1, and Bennett's bound is asymptotically 
    equivalent to Bernstein's. \\
    As $x \to \infty$, the taylor expansion of $h(x)$ gives $h(x) = x\log(x) - x + O(\log(x)) \sim x\log(x)$, so
    \[
        \frac{2(1 + x/3)}{x^2} h(x) = \frac{2(1 + x/3)}{x^2} (x\log(x)) \sim \frac{2}{3}\log(x) = \infty
    \]
    Thus, as $x \to \infty$, the ratio approaches infinity, and Bennett's bound is 
    asymptotically better than Bernstein's.
\end{enumerate}

\section{}
We are given the $\psi_\alpha$ Orlicz norm of a random variable $X$ defined as:
\[
    ||X||_{\psi_\alpha} = \inf\left\{K > 0 : E\left[\psi_\alpha \left(\frac{|X|}{K}\right)\right] \leq 1\right\}
\]
where $\alpha > 1$ and $\psi_\alpha(x) = e^{x^\alpha} - 1$. Using this, the
Orlicz norm can be rewritten as:
\[
    ||X||_{\psi_\alpha} = \inf\left\{K > 0 : E\left[e^{|X|^\alpha/K^\alpha}\right] \leq 2\right\}
\]
Suppose $||X||_{\psi_\alpha} < \infty$. By the definition of the infimum,
there exists a finite $K > 0$ where $E[e^{|X|^\alpha/K^\alpha}] \leq 2$. Then, 
for any $t > 0$, using Markov's inequality to the nonnegative random variable,
we get that $Y = e^{|X|^\alpha/K^\alpha}$. This gives:
\begin{align*}
    P(|X| \geq t) &= P(e^{|X|^\alpha/K^\alpha} \geq e^{t^\alpha/K^\alpha}) \\
    &\leq \frac{E[e^{|X|^\alpha/K^\alpha}]}{e^{t^\alpha/K^\alpha}} \\
    &\leq 2 e^{-t^\alpha/K^\alpha}
\end{align*}
Choosing $C = 2$ and $c = 1/K^\alpha$, we have that $P(|X| \geq t) \leq C e^{-ct^\alpha}$.
Thus, $||X||_{\psi_\alpha} < \infty$ implies that there exists a $C, c > 0$ such 
that $P(|X| \geq t) \leq C e^{-ct^\alpha}$. \\

\noindent
Now, we want to show that there exists a finite $K > 0$ such that $E[e^{|X|^\alpha/K^\alpha}] \leq 2$ 
if there exists a $C, c > 0$ such that $P(|X| \geq t) \leq C e^{-ct^\alpha}$. Let
$Y = |X|^\alpha \geq 0$, then for any $u \geq 0$:
\begin{align*}
    P(Y \geq u) &= P(|X|^\alpha \geq u) = P(|X| \geq u^{1/\alpha}) \\
    &\leq C e^{-c(u^{1/\alpha})} = C e^{-c u}
\end{align*}
For any $a > 0$, $e^{aY} = 1 + \int_0^Y a e^{au} du$. The tail integration identity
for nonnegative random variables gives us that $E[e^{aY}] = 1 + \int_0^\infty a e^{au} P(Y \geq u) du$. 
Using the tail bound on $P(Y \geq u)$, we have:
\begin{align*}
    E[e^{aY}] &\leq 1 + \int_0^\infty a e^{au} C e^{-cu} du \\
    &= 1 + aC \int_0^\infty e^{-(c - a)u} du \\
    &= 1 + \frac{aC}{c - a}
\end{align*}
Then if we choose $a = \frac{c}{C + 1}$, since $C, c > 0$, $0 < a < c$, and we have:
\begin{align*}
    E[e^{aY}] &\leq 1 + \frac{C(\frac{c}{C + 1})}{c - \frac{c}{C + 1}} \\
    &= 1 + \frac{\frac{Cc}{C + 1}}{\frac{cC}{C + 1}} \\
    &= 1 + 1 = 2
\end{align*}
Thus, $E[e^{aY}] \leq 2$. Recall that $Y = |X|^\alpha$, so we choose $K > 0$
such that $\frac{1}{K^\alpha} = a = \frac{c}{C + 1}$, and solving for $K$
gives us $K = \left(\frac{C + 1}{c}\right)^{1/\alpha}$. Then we have that
$E[e^{|X|^\alpha/K^\alpha}] \leq 2$, and then by the definition of the Orlicz norm,
there exists a finite $K > 0$. Because both implications hold, we have that:
\[
    ||X||_{\psi_\alpha} < \infty \iff \exists C, c > 0 : P(|X| \geq t) \leq C e^{-ct^\alpha} \quad \forall t \geq 0
\]

\noindent 
\textbf{AI Usage:} All homework problems were written in VSCode and compiled using LaTeX, with 
assistance from ChatGPT and Github Copilot inline suggestions.

\end{document}